Table~\ref{tab:distilled_perf} shows the performance and speed comparison between the Diffusion Decoder (50 steps) and Diffusion Decoder Turbo (8 steps) across five benchmarks. Speed is measured on a single GPU at $1024 \times 1024$ resolution with a batch size of 1, using BF16 precision.

\begin{table}[!htb]
\centering
\setlength\tabcolsep{10pt}
\renewcommand{\arraystretch}{1.0}
\caption{Performance and Speed Comparison between Diffusion Decoder and Diffusion Decoder Turbo.}
\resizebox{1.0\textwidth}{!}{%
\begin{tabular}{cc|ccccc}
\toprule
\textbf{Method} & \textbf{Speed (s / img)}
  & \textbf{GenEval} & \textbf{DPG} & \textbf{UniGenBench}
  & \textbf{OneIG-EN}
  & \textbf{WISE} 
  \\ 
\midrule
Diffusion Decoder (50 steps)
  & 32.95 & 0.89 & 87.76 & 79.63 & 0.505 & 0.68 \\
\midrule
Diffusion Decoder Turbo (8 steps)
  & 2.90 & 0.87 & 87.24 & 79.76 &0.500 & 0.68  \\
\bottomrule
\end{tabular}%
\label{tab:distilled_perf}
}
\end{table}

Through few-step acceleration, the Diffusion Decoder Turbo achieves an 11.4$\times$ speedup (from 32.95\,s/img to 2.90\,s/img) while maintaining competitive performance across all benchmarks.
%
Specifically, it retains a GenEval score of 0.87 (vs.\ 0.89 for the base) and a DPG score of 87.24 (vs.\ 87.76). 
%
Similarly, on the UniGenBench, OneIG-EN, and WISE benchmarks, Diffusion Decoder Turbo achieves performance on par with the original decoder.
%
Moreover, as shown in Figure~\ref{fig:decoder_comp}, the visual quality remains virtually indistinguishable.

\begin{figure}[!t]
    \centering
    \vspace{-0.1in}
    \includegraphics[width=1.0\linewidth]{figs/decoder_comp.pdf}
    \vspace{-0.23in}
    \caption{Visual Comparison of the Decoder and the Distilled Version.}
    \label{fig:decoder_comp}
    \vspace{-0.13in}
\end{figure}



